\subsection{\foreign{Linear Integer Arithmetic}}

Seperti pemetaan injektif \(\boolmap\), di sini didefinisikan pemetaan injektif \(\intmap\) dari \(V\) ke variabel \foreign{integer} di SMT. Dengan membatasi \(0 \leq \intmap(v_i) \leq 1\), seperti halnya koresponden keanggotaan simpul dalam \(W\) bergantung dari nilai \(\True\) untuk \(\boolmap(v_i)\), nilai \(1\) untuk \(\intmap(v_i)\) menunjukkan korespondensi keanggotaan dan nilai \(0\) untuk tidak. Sehingga, pemetaan \foreign{resolving set} ke dalam \foreign{constraint boolean satisfiability} dapat diangkat secara langsung ke bentuk \foreign{linear integer arithmetic} dengan melakukan translasi berikut:
\[
	\bigvee \limits_{\displaystyle v \in W} \boolmap(v)
	\iff
	\sum \limits_{\displaystyle v \in W} \intmap(v) \enspace \geq 1
\]

Dan didapat bentuk dari \autoref{eq:boolean-satisfiability-natural} dalam \foreign{linear integer arithmetic} sebagai berikut:
\begin{equation}
	\bigwedge \limits_{\displaystyle P_{*j} \in P} \enspace \sum \limits_{\displaystyle v \in V[\sim\!P_{*j}]} \intmap(v) \enspace \geq 1
	\label{eq:linear-integer-arithmetic-natural}
\end{equation}

\begin{example}
	Ambil graf \autoref{fig:algorithm-example}. Misalnya \(\intmap(v_i) = x_i\), maka \foreign{constraint} yang berkorespondensi dengan \foreign{resolving set}-nya dalam \foreign{linear integer arithmetic} adalah:
	\[
		\smash[b]{
		\begin{array}{c}
			x_1 + x_2 \geq 1 \\
			x_1 + x_3 \geq 1 \\
			x_2 + x_3 \geq 1 \\
			x_4 + x_5 \geq 1 \\ [8pt]
			0 \leq x_1 \leq 1,\enspace
			0 \leq x_2 \leq 1,\enspace
			0 \leq x_3 \leq 1,\enspace
			0 \leq x_4 \leq 1,\enspace
			0 \leq x_5 \leq 1
		\end{array}
		} \vspace{1em}
	\]
\end{example}

\clearpage

Selanjutnya, untuk \foreign{constraint} \(\AtLeast\) dan \(\AtMost\), dapat dinyatakan dalam \foreign{linear integer arithmetic} dengan lebih sederhana sebagai:
\begin{equation}
	\sum \limits_{\displaystyle v \in V} \intmap(v) \enspace = k
	\label{eq:linear-integer-arithmetic-cardinality}
\end{equation}

Perhatikan bentuk jumlahan ini yang lebih ekspresif dari pada \(\AtLeast\) dan \(\AtMost\). Menggunakan persamaan ini, \autoref{eq:linear-integer-arithmetic-natural} dapat ditransformasi menjadi relatif terhadap kardinalitas \(k\) seperti berikut:
\begin{equation}
	\bigwedge \limits_{\displaystyle P_{*j} \in P} \enspace \sum \limits_{\displaystyle v \in V[P_{*j}]} \intmap(v) \enspace \leq k - 1
	\label{eq:linear-integer-arithmetic}
\end{equation}

Sama seperti transformasi \autoref{eq:boolean-satisfiability-natural} ke \autoref{eq:boolean-satisfiability} yang membantu \foreign{conflict driven learning}, transformasi di \foreign{linear integer arithmetic} dari \autoref{eq:linear-integer-arithmetic-natural} ke \autoref{eq:linear-integer-arithmetic} juga membantu SMT dalam pemangkasan ruang pencarian.

\begin{example}
	Untuk graf \autoref{fig:algorithm-example}, \foreign{constraint} yang berkorespondensi dengan \foreign{resolving set} berkardinalitas \(k\) adalah:
	\[
		\smash[b]{
		\begin{array}{c}
			x_1 + x_2 + x_3 + x_4 + x_5 = k \\ [8pt]
			x_3 + x_4 + x_5 \leq k - 1 \\
			x_2 + x_4 + x_5 \leq k - 1 \\
			x_1 + x_4 + x_5 \leq k - 1 \\
			x_1 + x_2 + x_3 \leq k - 1 \\ [8pt]
			0 \leq x_1 \leq 1,\enspace
			0 \leq x_2 \leq 1,\enspace
			0 \leq x_3 \leq 1,\enspace
			0 \leq x_4 \leq 1,\enspace
			0 \leq x_5 \leq 1
		\end{array}
		} \vspace{1.5em}
	\]
	\label{ex:linear-integer-arithmetic-constraint}
\end{example}

Dengan demikian, didapatkan desain algoritma yang menggunakan \foreign{linear integer arithmetic}.

\begin{algorithm}
	\foreign{Input}:
	\begin{itemize}
		\item Representasi graf dalam representasi matriks ketetanggaan \(M\).
		\item Mencari dimensi metrik atau dimensi metrik sisi?
	\end{itemize}
	Prosedur:
	\begin{enumerate}
		\item[1--5.] Sama seperti pada \autoref{alg:boolean-satisfiability}.
		\setcounter{enumi}{5}
		\item Tambahkan \foreign{constraint} \(0 \leq \intmap(v_i) \leq 1\) untuk semua \(v_i \in V\) ke \foreign{solver} SMT.
		\item Mulai variabel \(k\) dari banyaknya simpul, yaitu \(k := |V|\).
		\item Selama \(k \geq 0\):
			\begin{enumerate}
					\item Tambahkan \foreign{constraint} berdasarkan \autoref{eq:linear-integer-arithmetic-cardinality}, dan berdasarkan \autoref{eq:linear-integer-arithmetic} dari \(P\), ke \foreign{solver} SMT.
					\item Jalankan \foreign{solver} untuk mencari \foreign{model}.
					\item Jika dihasilkan \(\unsat\), lompat ke tahap (10).
					\item Hapus \foreign{constraint} yang ditambahkan dari tahap (8a) sebelumnya.
					\item Kurangi \(k\) sebanyak satu, yaitu \(k := k - 1\).
			\end{enumerate}
		\item Simpulkan bahwa 0 adalah dimensi metrik atau dimensi metrik sisi, dan \texttt{BERHENTI}.
		\item Simpulkan bahwa \(k + 1\) adalah dimensi metrik atau dimensi metrik sisi.
	\end{enumerate}
	\caption{\foreign{Linear Integer Arithmetic}}
	\label{alg:linear-integer-arithmetic}
\end{algorithm}

\begin{example}
	Ambil graf \autoref{fig:algorithm-example}, dan cari dimensi metriknya menggunakan \autoref{alg:linear-integer-arithmetic}. Setelah menetapkan \foreign{constraint} \(0 \leq x_i \leq 1\) untuk setiap variabel \foreign{integer}-nya, \foreign{constraint} yang digunakan pada setiap tahapannya bersesuaian dengan yang telah tertera pada \autoref{ex:linear-integer-arithmetic-constraint}.
	\begin{itemize}
		\item Untuk \(k = 5\), ditemukan \(\sat\).
		\item Untuk \(k = 4\), ditemukan \(\sat\).
		\item Untuk \(k = 3\), ditemukan \(\sat\).
		\item Untuk \(k = 2\), ditemukan \(\unsat\).
	\end{itemize}
	Disimpulkan bahwa \(2 + 1 = 3\) adalah dimensi metrik dari graf tersebut.
\end{example}
